\documentclass[11pt]{article}
\usepackage[margin=1in]{geometry}
\usepackage{amsmath,amssymb,amsthm}
\usepackage[hidelinks]{hyperref}
\newtheorem{theorem}{Theorem}
\newcommand{\poly}{\mathrm{poly}}
\begin{document}
\begin{center}
{\LARGE Zhang's uncountable triple rainbow problem}
\end{center}
\vspace{0.6em}
\begin{abstract}
We give a ZFC counterexample to both clauses of Question 5.5 of
Jing Zhang. The construction is a matching rule indexed by the two
largest vertices of a triple, using a strong colouring of pairs.
\end{abstract}

For a set \(A\) of ordinals, a colouring of \([A]^r\) is
\(2\)-bounded if every colour class has at most two members. A subset
of \(A\) is a \emph{rainbow} if the colouring is injective on its
\(r\)-element subsets.

\begin{theorem}
There is a \(2\)-bounded \(f:[\omega_1]^3\to\omega_1\) such that,
for every \(u<v<\omega_1\) and every uncountable
\(X\subseteq\omega_1\setminus(v+1)\), some \(a<b\) in \(X\) satisfy
\[
   f(\{u,a,b\})=f(\{v,a,b\}).
\]
In particular, \(f\) has no uncountable rainbow.
\end{theorem}

\begin{proof}
By Todor\v{c}evi\'c's square-bracket theorem \cite{tod}, take
\(c:[\omega_1]^2\to\omega_1\) taking all values on the pairs of every
uncountable set. Fix an enumeration
\(\langle p_\xi:\xi<\omega_1\rangle\) of the pairs of ordinals.
For \(a<b\), the single pair \(p_{c(\{a,b\})}\) instructs the
matching at \(\{a,b\}\). If that pair is \(\{u,v\}\subseteq a\),
give the two triples \(\{u,a,b\}\) and \(\{v,a,b\}\) one new common
colour. Use distinct such colours at distinct top pairs, and a new
singleton colour for each remaining triple. A triple has a unique
top pair, so these instructions are disjoint and the resulting
colouring is \(2\)-bounded, using at most \(\aleph_1\) colours.

For \(u,v,X\) as in the statement, choose \(\xi\) with
\(p_\xi=\{u,v\}\). Some \(a<b\) in \(X\) satisfy \(c(\{a,b\})=\xi\).
Both \(u,v\) lie below \(a\), so this instruction gives the claimed
collision. For any uncountable candidate rainbow, take \(u,v\) inside
it and let \(X\) be its tail above them.
\end{proof}

Thus the proposed relation in \cite[Question 5.5]{zhang} is false in
ZFC, and in particular PFA cannot imply it. Appending a common tuple
of larger vertices to each colliding triple, and recording that tuple
in the colour, gives the corresponding result for every finite
arity at least three. This does not decide the separate question
about arbitrary countable order types.

\begin{thebibliography}{9}
\bibitem{tod}
S. Todor\v{c}evi\'c, \emph{Partitioning pairs of countable ordinals},
Acta Math. 159 (1987), 261--294.
\href{https://doi.org/10.1007/BF02392561}{doi:10.1007/BF02392561}.

\bibitem{zhang}
J. Zhang, \emph{Some remarks on uncountable rainbow Ramsey theory},
Proc. Amer. Math. Soc. 151 (2023), 865--880.
\href{https://arxiv.org/abs/1809.00649}{arXiv:1809.00649}.
\end{thebibliography}
\end{document}
